\documentclass[11pt]{article}
\usepackage[margin=2.1cm]{geometry}
\usepackage{booktabs}
\usepackage{amsmath}
\usepackage{amssymb}
\usepackage{enumitem}
\usepackage{xcolor}
\usepackage[T1]{fontenc}
\usepackage{lmodern}

\setlength{\parindent}{0pt}
\setlength{\parskip}{3.6pt}
\newcommand{\code}[1]{\texttt{#1}}
\newcommand{\hr}{\vspace{2pt}\hrule\vspace{6pt}}

\begin{document}

\begin{center}
{\Large\bfseries EE\,362 --- Computer Organization Laboratory}\\[2pt]
{\large\bfseries Lab 0 --- VHDL Fundamentals}\\[4pt]
Fall 2026 \quad$\cdot$\quad Lab TA: Gandhi Torres \quad$\cdot$\quad Sections: Mon / Tue, 3:45--5:45~PM, ECIII 237
\end{center}
\hr

\section*{1. Objective}

Learn the shape of a VHDL design --- the \code{entity} that declares an interface
and the \code{architecture} that implements it --- and verify four small
combinational circuits by simulation rather than by inspection.

Every one of these four circuits reappears later in the course. Two of them end up
inside your processor.

\section*{2. The method, which does not change all semester}

\begin{quote}
\textbf{specification $\rightarrow$ truth table $\rightarrow$ Boolean equations
$\rightarrow$ VHDL $\rightarrow$ simulation $\rightarrow$ back to the table}
\end{quote}

Today it is a four-row table and one equation. In week twelve it is an eight-row
table with sixteen output columns, and nothing else will have changed.

\subsection*{How to write a truth table}

\begin{itemize}[leftmargin=1.6em,itemsep=1pt]
\item \textbf{Number every row.} Your report says ``row 5'', never a paragraph
describing row 5.
\item Inputs on the left, outputs on the right, separated by a visible bar.
\item \code{--} means \emph{don't care}.
\item With $n$ binary inputs an exhaustive table has $2^n$ rows. Four inputs or
fewer: enumerate them all.
\end{itemize}

\subsection*{How to write the derivation}

Three lines, every time:

\begin{align*}
\text{(1) canonical} \quad & \text{sum of minterms, straight off the table} \\
\text{(2) simplified} \quad & \text{\textbf{and name the law or K-map grouping used}} \\
\text{(3) VHDL} \quad & \text{the line of code that implements it}
\end{align*}

``It simplified'' is not line (2). Name what you did.

\section*{3. The four circuits}

\subsection*{Exercise 1 --- majority of three}

$Y = 1$ when at least two of $A$, $B$, $C$ are $1$.

Three inputs, so an exhaustive table has 8 rows. Fill it in, write
$Y = \sum m(\ldots)$, then simplify to a three-term expression and name the law.

\subsection*{Exercise 2 --- 2-to-1 multiplexer}

$Y = D_0$ when $S = 0$, $Y = D_1$ when $S = 1$.

Write it \emph{both} ways: the exhaustive 8-row table, and the 2-row table that
names a \emph{source} rather than a value. You will use the second form for the
rest of the course, and Lab 3 explains why.

\subsection*{Exercise 3 --- 2-to-4 decoder with enable}

Exactly one output is high when $E = 1$; all outputs low when $E = 0$.

Write one equation per output. Notice what they have in common --- that pattern
has a name, and you meet it again in Lab 6.

\textbf{Warning:} the behavioural example given in class has a bug. Find it. It is
Exercise 3's real lesson.

\subsection*{Exercise 4 --- full adder}

$A + B + C_{in}$, producing $Sum$ and $C_{out}$.

Two outputs, so two equations. One of them does \emph{not} simplify by K-map ---
work out why, and say what operator it needs instead. The other one you have
already derived today; say which exercise it was.

\section*{4. What to build}

For each of the four circuits:

\begin{enumerate}[leftmargin=2em,itemsep=2pt]
\item the truth table, rows numbered;
\item the three-line derivation, with the law named;
\item a VHDL \code{entity} and \code{architecture};
\item a \textbf{self-checking} testbench --- \code{assert} statements, not
eyeballing a waveform;
\item the waveform, and the transcript line showing the checks passed.
\end{enumerate}

For Exercises 1 and 4 write \textbf{two architectures} of the same entity ---
\code{behavioral} and \code{structural} --- and have one testbench exercise both,
asserting they agree on every vector.

\section*{5. What to enter in Quartus}

You build this project by hand. It is the only time you will --- from Lab 1 the
\code{.qpf} is supplied --- and doing it once is why.

\begin{enumerate}[leftmargin=2em,itemsep=2pt]
\item \code{File $\rightarrow$ New Project Wizard}
\item \textbf{Working directory:} a short local path, \emph{no spaces}, not inside
OneDrive. \code{C:\textbackslash ee362\textbackslash lab00} is fine.
\item \textbf{Project name} and \textbf{top-level entity name}: these must match
the \code{entity} name in your VHDL file exactly.
\item \textbf{Family:} Cyclone V. \quad \textbf{Device:} \code{5CSEMA4U23C6}.
\item \textbf{EDA Tool Settings $\rightarrow$ Simulation:} tool
\code{Questa Intel FPGA (VHDL)}, format VHDL.
\item Add your \code{.vhd} files: \code{Project $\rightarrow$ Add/Remove Files}.
\item \code{Assignments $\rightarrow$ Settings $\rightarrow$ EDA Tool Settings
$\rightarrow$ Simulation $\rightarrow$ Compile test bench} $\rightarrow$
\code{Test Benches...} $\rightarrow$ \code{New}, and give it your testbench entity
name for both the test bench name and the top-level module.
\end{enumerate}

Then the flow:

\begin{center}
\begin{tabular}{lll}
\toprule
& what & why \\
\midrule
1 & \code{Processing $\rightarrow$ Start $\rightarrow$ Analysis \& Elaboration} & fast syntax check \\
2 & \code{Tools $\rightarrow$ Netlist Viewers $\rightarrow$ RTL Viewer} & did I describe what I meant? \\
3 & \code{Processing $\rightarrow$ Start Compilation} (Ctrl+L) & full build \\
4 & Compilation Report & errors, warnings, resources \\
5 & \code{Tools $\rightarrow$ Run Simulation Tool $\rightarrow$ RTL Simulation} & does it behave? \\
\bottomrule
\end{tabular}
\end{center}

\textbf{Do step 2 every time.} It costs ten seconds and catches most design
mistakes before you have spent a compile on them.

\section*{6. Report --- \code{lab00.tex}}

All seven sections, for all four circuits. Specifically: the four truth tables
with numbered rows, the four derivations with laws named, the code included by
\code{\textbackslash lstinputlisting} (never pasted), the coverage argument for
each testbench, and a Discussion describing one real bug you hit.

\vspace{10pt}
\hr
{\small Two of these four circuits end up inside your processor, and a third is
the mechanism behind the register file. Keep your files.}

\end{document}
